\subsection*{物体没有换，描述它的语言换了}
两个人看同一箭头，一个用水平与竖直方向报坐标，一个用两根斜基报坐标。数字不同不代表向量不同。先约定 $P$ 接收新坐标、输出旧坐标。因为新坐标 $(1,0)$ 就是第一根新基，$P$ 的第一列必须是这根基在旧坐标下的描述。

已知动作 $A$ 使用旧坐标，却收到新坐标输入，就要先翻译、再执行、再翻译回来。实际顺序是 $P,A,P^{-1}$，写成乘积是 $P^{-1}AP$。只要记清三句指令，逆矩阵的位置不必靠猜。

\subsection*{好语言让同一动作变简单}
交换横纵坐标的动作，把 $(1,1)$ 留在原处，把 $(1,-1)$ 反向。选它们为新基，动作就从“交换两个坐标”变为“保留第一分量、取反第二分量”，因此出现对角表示。

例如 $(3,1)$ 的新坐标由 $a+b=3,a-b=1$ 得到 $(2,1)$。新矩阵把它变成 $(2,-1)$，再翻译回去得到 $(1,3)$，与原动作相符。走完这三步，往往比直接展开相似变换更容易发现方向约定错误。

换基改变描述，主动变换改变对象；两种解释共享矩阵形式，很容易混淆。阅读推导时要明确当前向量是物体、旧坐标还是新坐标。下一节寻找特征向量，就是寻找一套好语言，让重复作用变成独立方向上的简单伸缩。
